\documentclass{beamer}

\usepackage[spanish,activeacute,es-tabla]{babel}
\usepackage[utf8]{inputenc} % mac utf8
\usepackage{cancel}
\usepackage{booktabs}


\usepackage{graphicx}
\DeclareGraphicsExtensions{.png,.pdf,.jpg}


\usefonttheme{professionalfonts} % fuentes de LaTeX 
\usetheme{Boadilla}	% tema escogido en este ejemplo 


\setbeamertemplate{caption}[numbered]
\newtheorem{Teorema}{Teorema} 
\newtheorem{Definicion}{Definici\’on} 
\newtheorem{Prueba}{Prueba}

\theoremstyle{plain}

\newtheorem{teorema} [Teorema] {Teorema}
\newtheorem{definicion} [Definicion]{Definición} 
\newtheorem{lema} [theorem] {Lema}


\theoremstyle{example}
\newtheorem{Ejercicio}{Ejercicio} 
\newtheorem{ejercicio} [Ejercicio]{Ejercicio} 

\theoremstyle{example}
\newtheorem{Ejemplo}{Ejemplo} 
\newtheorem{ejemplo} [Ejemplo] {Ejemplo}

\setbeamertemplate{theorems}[numbered] 

% 3) Colores por defecto del style 'example' (fallback)
\setbeamercolor{block title example}{bg=green!40!black, fg=white}
\setbeamercolor{block body example}{bg=green!10, fg=black}

% 4) Colores por ENTORNO (locales) usando etoolbox
\usepackage{etoolbox}

% -- 'Ejemplo' en naranja
\AtBeginEnvironment{ejemplo}{%
  \setbeamercolor{block title example}{bg=orange!60!black, fg=white}%
  \setbeamercolor{block body example}{bg=orange!10, fg=black}%
}
% -- 'Ejercicio' en violeta
\AtBeginEnvironment{ejercicio}{%
  \setbeamercolor{block title example}{bg=violet!60!black, fg=white}%
  \setbeamercolor{block body example}{bg=violet!10, fg=black}%
}





\title[]{Tema 7. Teoría Codificación Algebraica}
\author[José A Montenegro (jmmontes@uma.es)]{José A. Montenegro}
\institute{Dpto. Lenguajes y Ciencias de la Computación\\
ETSI Informática. Universidad de Málaga\\
jmmontes@uma.es
%\href{https://twitter.com//monteDocencia}{
%
\includegraphics[height=0.03\textheight]{twitter.jpg}}
}
\date{\today}

%###########################################################################
% FootLine
%###########################################################################

\setbeamertemplate{footline} {%
\leavevmode%
  \hbox{\begin{beamercolorbox}[wd=.5\paperwidth,ht=2.5ex,dp=1.125ex,leftskip=.3cm plus1fill,rightskip=.3cm]{author in head/foot}%
    \usebeamerfont{author in head/foot}\insertshortauthor 
  \end{beamercolorbox}%
  
  \begin{beamercolorbox}[wd=.5\paperwidth,ht=2.5ex,dp=1.125ex,leftskip=.3cm,rightskip=.3cm plus1fil]{title in head/foot}%
    Teoría de la Información y Codificación \hfill    \insertpagenumber /\insertdocumentendpage
  \end{beamercolorbox}}%
  \vskip0pt%
}%

%###########################################################################
% Document
%###########################################################################

\begin{document}

%\setbeamertemplate{headline}[title slide]
%\setbeamertemplate{footline}[title slide]


%###########################################################################
% Tíitulo
%###########################################################################
\begin{frame}

\includegraphics[height=0.15\textheight]{logoUma.jpg} \hspace*{3.5cm}

\includegraphics[height=0.17\textheight]{logoInf.jpg}   \hspace*{3.5cm}

\includegraphics[height=0.17\textheight]{logoLcc.jpg}  

\titlepage


\setbeamertemplate{footline}[body]
\setbeamertemplate{headline}[body]
\end{frame}

%###########################################################################
% Contenidos
%###########################################################################

\section[Resumen]{}
\frame{\tableofcontents}

%%%%%%%%%%%%%%%%%%%%%%%%%%%%%%%%
%%%%%%%%%%%%%%%%%%%%%%%%%%%%%%%%
%SECCION %%%%%%%%%%%%%%%%%%%%%%%%%%%%%%%
%%%%%%%%%%%%%%%%%%%%%%%%%%%%%%%%
%%%%%%%%%%%%%%%%%%%%%% 

\begin{frame}   \footnotesize
{
\setbeamercolor{block title}{bg=blue!60!black, fg=white}
\setbeamercolor{block body}{bg=blue!10, fg=black}
\begin{block}{Sección 1}
Códigos Hamming
\end{block}
}
\end{frame} 

\section{Códigos Hamming}

%%%%%%%%%%%%%%%%%%%%%%%%%%%%%%%%
\begin{frame}  \frametitle{Códigos Hamming}\footnotesize

\begin{itemize} 

\item  En el capitulo anterior mostramos que los métodos algebraicos pueden ser utilizados para construir códigos útiles. Nos centramos en una importante \textbf{familia de códigos}, descubiertos en 1950 por R.W. \textbf{Hamming}.

\item Acorde con el \textbf{teorema 1 del tema 6} un código binario lineal con $\delta \geq 3$ puede ser definido por una \textbf{matriz de verificación} en la cual todas las columnas son distintas y no son igual a cero.

\item Si tenemos m filas, entonces hay $2^m$ vectores columnas de longitud m, por lo que el máximo número de columnas distintas no ceros es $2^m -1$.

\item El orden en cual las columnas son listadas no es relevante, aunque en la practica establecer un orden puede resultar útil.

%\item Reordenar las filas solamente implica una reordenación de los bits de cada palabra codificada, y no afectan a los parámetros del código. 

\item Dos códigos que difieren en el orden de las filas son equivalentes.
\end{itemize} 

\footnotesize

\begin{definicion}[Código de Hamming]
Un código definido por una matriz de verificación binaria con m filas y $2^m -1$ columnas distintas y ninguna  cero es un \textbf{código binario de Hamming}.
\end{definicion}

\end{frame} 

%%%%%%%%%%%%%%%%%%%%%%%%%%%%%%%%
\begin{frame}  \small

\begin{ejemplo}
Establezca la matriz de verificación y los parámetros del código Hamming con m = 3.\\
\end{ejemplo}
\textit{\underline{Solución}}:\\
\vspace{0.3cm}
Tenemos $2^3 -1=7$ columnas. \\
\vspace*{5pt}
Una forma de listarlas todas sería establecer un orden de diccionario.\\
\vspace*{5pt}
 Pero, como observaremos posteriormente, hay razones para utilizar el orden numérico, el cual es el orden correspondiente a la representación binarias de los números 1, 2, 3, 4, 5, 6, 7:

\[
    H_1=   \left( \begin{array}{ccccccc}
     				  0 &   0   &   0  &   1 &	1 & 1 & 1\\ 
					  0 &   1   &   1  &   0 &	0 & 1 & 1\\ 
					  1 &   0   &   1  &   0 &	1 & 0 & 1
           			\end{array}
    		\right) 	    		
\]



\end{frame} 

%%%%%%%%%%%%%%%%%%%%%%%%%%%%%%%%
\begin{frame}  \small

\textit{\underline{Solución}}:\\
\vspace{0.3cm}

De forma alternativa podemos ordenar las columnas de forma que aquellas con peso 1  se sitúan al final, obteniendo una matriz de verificación en la \textit{forma estándar}:

\[
    H_2=   \left( \begin{array}{ccccccc}
     				  0 &   1   &   1  &   1 &	1 & 0 & 0\\ 
					  1 &   0   &   1  &   1 &	0 & 1 & 0\\ 
					  1 &   1   &   0  &   1 &	0 & 0 & 1
           			\end{array}
    		\right) 	    		
\]

Los códigos definidos por las matrices $H_1$ y $H_2$ son equivalentes.\\
\vspace*{5pt}
Los parámetros $(n,k,\delta)$ se pueden establecer fácilmente.\\
\vspace*{5pt}
La longitud de palabra es n=7 y la dimensión es k = n - m = 4. \\
\vspace*{5pt}
La distancia mínima $\delta$ es al menos 3, ya que todas las columnas de $H_1$ son distintas y no ceros, es fácil verificar que $c_1$=1110000 ($c_1 H_2=0$) es una palabra codificada con peso 3, por lo que $\delta =3$.


\end{frame} 
%%%%%%%%%%%%%%%%%%%%%%%%%%%%%%%%
\begin{frame}  \footnotesize  


\begin{itemize}

 \item Generalmente, para cada $m\geq 3$, hay un \textbf{código de Hamming binario} con parámetros
\begin{center}
 $n=2^m-1$, $k=2^m-1-m$ o $k=n-m$ y $\delta=3$.
\end{center}

\begin{teorema}\label{teo:CBH}
En un \textbf{código binario de Hamming}, todas la palabras en $\mathbb{F}_2^n$ son una \textbf{palabra codificada} o están a una \textbf{distancia 1 de una palabra codificada}.
\end{teorema}

\item  Generalmente, un código corrector r tiene la propiedad que $\delta\geq 2r+1$, por lo que los \textbf{vecindarios} $N_r(c)$ ($c \in C)$ son mutuamente disjuntos.

\item Sin embargo hay  ``\textbf{huecos}'' que no cubren esos vecindarios, por lo que algunas palabras no están en ninguno de esos vecindarios $N_r(c)$.

\end{itemize} 

\begin{definicion}[Código Perfecto]

Un código corrector r C es un código perfecto si todas las palabras están en uno de los vecindarios $N_r(c)$ para algún $c \in C$.
\end{definicion}


\begin{itemize} 

\item El teorema \ref{teo:CBH} establece que un \textbf{código binario Hamming} es un \textbf{código perfecto con r=1}. 


\end{itemize} \end{frame} 





%%%%%%%%%%%%%%%%%%%%%%%%%%%%%%%%
\begin{frame}  \small  

\vspace{0.3cm}
\begin{itemize}
    \item Con \textbf{m} = 3, \textbf{n} sería $2^3-1$ = 7. \textbf{k} = 7-3 = 4.

    \item Tenemos $2^\mathbf{k}=2^4=$ \textbf{16} palabras codificadas c y cada vecindario $\mathbf{N_i(c)}$ tiene 1+7 = 1+n = \textbf{8} palabras. 

    \item Estos \textbf{16 vecindarios} son disjuntos y ya que 16 x 8= \textbf{128} = $2^n = 2^7$ es el número total de palabras, el código es perfecto (véase siguiente Figura).

\end{itemize}

\begin{figure}[!hbtp]

\begin{center}
%\includegraphics[height=0.5\textheight]{PerfectCode.pdf} 
\includegraphics[height=0.5\textheight]{Tema7/hamming7_neighborhoods.png} 
\end{center}
\caption{El código Hamming con parámetros  (n=7, k=4, $\delta$= 3) es perfecto}\label{fig:perfectcode}
\end{figure}

\end{frame} 



%%%%%%%%%%%%%%%%%%%%%%%%%%%%%%%%
\begin{frame}  \small  
La mayor razón para utilizar métodos algebraicos es que las estructuras algebraicas pueden proporcionar buenos algoritmos para codificar y decodificar.\\ 
\vspace*{5pt}
Para los códigos binarios de Hamming, podemos utilizar las siguientes técnicas.

\begin{description}

\item[Codificación] 	Las columnas de la matriz de verificación H están ordenadas en \textbf{orden numérico}.\\
\vspace*{5pt}
Entonces los bits $x_1,x_2,x_4,\ldots,x_{2^{m-1}}$	son \textbf{bits de verificación}, y podemos escribir las ecuaciones que los definen en términos de los bits restantes (\textbf{mensaje}).\\
\vspace*{5pt}
En el caso de m=3, las ecuaciones derivadas de $H_1$ son:

\begin{center}
$x_1 =x_3 +x_5 +x_7$,	\\$x_2 =x_3 +x_6 +x_7$,\\	$x_4 =x_5 +x_6 +x_7$.
\end{center}



\end{description}
\end{frame} 
%%%%%%%%%%%%%%%%%%%%%%%%%%%%%%%%

\begin{frame}  \small


\begin{description}

\item[cont Codificación]

A veces es conveniente utilizar los primeros \textbf{k bits como los bits de mensaje} y el resto \textbf{n-k bits como bits de verificación}. \\
\vspace*{5pt}
Esto corresponde a escribir la matriz de verificación en su forma estándar $H_2$, las columnas son obtenidas de $H_1$ mediante permutación

\begin{center}
$1\mapsto 7, 2 \mapsto 6, 3 \mapsto 1, 4 \mapsto 5, 5\mapsto 2, 6  \mapsto 3, 7  \mapsto 4.$
\end{center}

Aplicando estas permutaciones y reordenando las ecuaciones, los bits de verificación son dados por:

\begin{center}
$x_5 =x_2 +x_3 +x_4$,\\	$x_6 =x_1 +x_3 +x_4$,\\	$x_7 =x_1 +x_2 +x_4.$
\end{center}


\end{description}
\end{frame} 

%%%%%%%%%%%%%%%%%%%%%%%%%%%%%%%%
%%%%%%%%%%%%%%%%%%%%%%%%%%%%%%%%
\begin{frame}  \small  

\begin{description}



\item[Decodificación] Cuando la \textbf{matriz de verificación} es ordenada en un \textbf{orden numérico}, la \textbf{regla de decodificación del síndrome} toma una forma muy simple.\\
\vspace*{5pt}
Hay n + 1 clases de equivalencia 

\begin{center}
$0+C$,	$e_1 +C$,	$e_2 +C$,	$\ldots$ ,	$e_n +C$,
\end{center}

donde $e_i$ denota la palabra $0\ldots 010\ldots 0$ en el cual el $i^{th}$ bit es 1. \\
\vspace*{10pt}
El \textbf{síndrome} de las clase de equivalencia $e_i + C$ es $He'_i$ el cual es igual a la $i^{th}$ columna de $H$. \\
\vspace*{15pt}
\textbf{La regla para corregir un error} es simple: el síndrome de la palabra recibida \textbf{z} es la \textbf{representación binaria} de $\mathbf{i}$, entonces el $\mathbf{i^{th}}$ \textcolor{red}{bit es erróneo}. 

\end{description}
\end{frame} 

%%%%%%%%%%%%%%%%%%%%%%%%%%%%%%%%
%%%%%%%%%%%%%%%%%%%%%%%%%%%%%%%%
\begin{frame}  \small \begin{ejercicio} 

Establece la matriz de verificación para el código Hamming de longitud 15, utilizando el orden numérico de las columnas. \\
\vspace*{5pt}
¿Cuantas palabras codificadas tenemos? \\
\vspace*{5pt}
¿Cuales de las siguientes palabras son palabras codificadas? 011010110110000, 100000100000011, 110010110111111, 011001110011001. \\
\vspace*{5pt}
Corrige aquellas que no son palabras codificadas, asumiendo que solamente puede ocurrir un error.\\
\vspace*{5pt}

\end{ejercicio} 
\end{frame} 

\begin{frame}  \small

\textit{\underline{Solución}}:\\ 
\vspace{0.3cm}
Los parámetros serían n=15, $n=2^m-1$; m = 4. k = n-m= 11.\\
\vspace{0.15cm}
Con lo cual tendremos $2^{k} = 2^{11}=2048$ \textbf{palabras codificadas}. 
\vspace{0.15cm}

La matriz es: 


\[
    H=   \left( \begin{array}{ccccccccccccccc}
0&	0	&0	&0&	0	&0	&0	&1	&1	&1	&1	&1	&1	&1	&1\\		
0&	0	&0	&1&	1	&1	&1	&0	&0	&0	&0	&1	&1	&1	&1\\
0&	1	&1	&0&	0	&1	&1	&0	&0	&1	&1	&0	&0	&1	&1\\
1&	0	&1	&0&	1	&0	&1	&0	&1	&0	&1	&0	&1	&0	&1
           			\end{array}
    		\right) 	    		
\]

\vspace{0.15cm}
$z_0 (011010110110000)$; H $z_0$ = 1010 = 10 = 011010110010000\\
$z_1 (100000100000011)$; H $z_1$ = 0111 = 7   = 100000000000011\\
$z_2 (110010110111111)$; H $z_2$ = 1000 = 8   = 110010100111111\\
$z_3 (011001110011001)$; H $z_3$ = 1000 = 0   = \textbf{palabra codificada}\\
\end{frame} 


\begin{frame}[fragile]  {hamming.py}\footnotesize  

\begin{verbatim}  
Usage: python3 hamming.py <m> [--encode <kbits> | --random] [--seed <int>]
[--flip <pos> | --flip-rand [count]] [--decode <nbits>] 
[--selftest [--tries N]]
\end{verbatim}
\tiny
\begin{verbatim}  
python3 hamming.py 4 --random --flip 2
m = 4
n = 2^4 - 1 = 15
k = n - m = 11

H (columnas 1..n):
r 1:  0 0 0 0 0 0 0 1 1 1 1 1 1 1 1
r 2:  0 0 0 1 1 1 1 0 0 0 0 1 1 1 1
r 3:  0 1 1 0 0 1 1 0 0 1 1 0 0 1 1
r 4:  1 0 1 0 1 0 1 0 1 0 1 0 1 0 1

== Encoding input (random) ==
data (11 bits):   00111011011
code (15 bits):   010101111011011
* Flipped positions: [2]
code with errors: 000101111011011

== Decoding ==
syndrome: 0100  (LSB-first)
single-bit error at position 2 (corrected)
code with errors: 000101111011011
corrected code:   010101111011011
recovered data (11 bits): 00111011011
\end{verbatim}

\end{frame} 


\begin{frame}[fragile]  {hamming\_neighborhoods.py}\footnotesize  

\begin{verbatim}  
usage: hamming_neighborhoods.py [-h] --m M [--png PNG] [--csv CSV] 
[--sample SAMPLE] [--seed SEED] [--codeword CODEWORD] [--all] [--no-plot]
\end{verbatim}
\tiny
\begin{verbatim}  
python3 hamming_neighborhoods.py --m 3
m=3  n=2^3-1=7  k=n-m=4  codewords=2^4=16
[CSV] guardado en: neighborhoods.csv  (vecindarios: 16)
[PNG] guardado en: neighborhoods.png  (vecindarios dibujados: 16)

[INFO] Se han generado TODOS los vecindarios (16).

Resumen (hasta 5):
  0000000 | size=8 | members=0000000, 0000001, 0000010, 0000100, 0001000, 0010000, 0100000, 1000000
  0001111 | size=8 | members=0000111, 0001011, 0001101, 0001110, 0001111, 0011111, 0101111, 1001111
  0010110 | size=8 | members=0000110, 0010010, 0010100, 0010110, 0010111, 0011110, 0110110, 1010110
  0011001 | size=8 | members=0001001, 0010001, 0011000, 0011001, 0011011, 0011101, 0111001, 1011001
  0100101 | size=8 | members=0000101, 0100001, 0100100, 0100101, 0100111, 0101101, 0110101, 1100101
... (ver CSV para el listado completo)
\end{verbatim}

\end{frame} 

%%%%%%%%%%%%%%%%%%%%%%%%%%%%%%%%
%%%%%%%%%%%%%%%%%%%%%%%%%%%%%%%%
%EJERCICIO %%%%%%%%%%%%%%%%%%%%%%%%%%%%%%%
%%%%%%%%%%%%%%%%%%%%%%%%%%%%%%%%



%9.2. Show that if the check matrix for a Hamming code is arranged so that the columns are in numerical order, then the word 11100 · · ·00 is a codeword.

%9.3. Prove that the number of cosets of the the binary Hamming code C with words of length n is n + 1, and show that a complete set of coset representatives is the set of n + 1 words 0, e1, e2, . . . en, as defined above.

%%%%%%%%%%%%%%%%%%%%%%%%%%%%%%%%
%\begin{frame}  \small  
%\begin{itemize} \item  \item \item \item 
%\end{itemize} \end{frame} 


%%%%%%%%%%%%%%%%%%%%%%%%%%%%%%%%
%\begin{frame}  \small  
%\end{frame} 
%%%%%%%%%%%%%%%%%%%%%%%%%%%%%%%%
%\begin{frame}  \small \begin{ejercicio} \end{ejercicio}\end{frame} 



%%%%%%%%%%%%%%%%%%%%%%%%%%%%%%%%
%%%%%%%%%%%%%%%%%%%%%%%%%%%%%%%%
%SECCION %%%%%%%%%%%%%%%%%%%%%%%%%%%%%%%
%%%%%%%%%%%%%%%%%%%%%%%%%%%%%%%%
%%%%%%%%%%%%%%%%%%%%%% 

\begin{frame}   \footnotesize
{
\setbeamercolor{block title}{bg=blue!60!black, fg=white}
\setbeamercolor{block body}{bg=blue!10, fg=black}
\begin{block}{Sección 2}
Códigos Cíclicos
\end{block}
}
\end{frame} 
\section{Códigos Cíclicos}


%%%%%%%%%%%%%%%%%%%%%%%%%%%%%%%%
\begin{frame}  \frametitle{Códigos Cíclicos}\footnotesize

\begin{itemize} 

\item  Los \textbf{códigos Hamming} proporciona una familia de códigos en los cuales la \textbf{dimensión del código} es $k$, y por tanto $|C| = 2^k$, puede ser tan largos como necesitemos, sólo tenemos que incrementar m, $n=2^m-1, k=n-m$, mientras que la \textbf{distancia mínima es constante}, $\delta$=3.

\begin{center}
 \textcolor{red}{¿Existe alguna familia de códigos que pueda crecer en longitud y además mejorar su capacidad de corrección de errores ($\delta$ crezca junto con n)?}
\end{center}

\item  Utilizaremos $a= a_0a_1\ldots a_{n-1}$ para representar una palabra en $\mathbb{F}_2^n$.

\item El \textbf{desplazamiento cíclico} de $a$ es una palabra: $\hat{a}=a_{n-1}$ $a_0$ $a_1$ $\ldots$ $a_{n-2}.$
\end{itemize} 

\begin{definicion}[Código Cíclico]

Un Conjunto $C\subseteq \mathbb{F}_2^n$ es un \textbf{código cíclico} si
\begin{enumerate}
\item es un código lineal,
\item $\hat{c}$ está en C, si c está en $C$, esto es: $c_0\,c_1\,c_2 \ldots c_{n-1} \in C	\Rightarrow c_{n-1}\,c_0\,c_1 \ldots c_{n-2} \in C$.
\end{enumerate}

La definición implica que si \textbf{C es cíclica} y $c\in C$ entonces la palabra: $c_i\,c_{i+1} \ldots c_{n-1}$ $\,c_0 \ldots c_{i-1}$ es obtenida de c mediante un \textbf{número de desplazamientos} y también está en C.
\end{definicion}

\end{frame} 

%%%%%%%%%%%%%%%%%%%%%%%%%%%%%%%%
\begin{frame}  \small  

\begin{ejemplo}
¿Cual de los siguientes códigos son cíclicos? 

\begin{center}
$C_1 =\{000,100,010,001\}, \;\;	C_2 =\{0000,1010,0101,1111\}$
\end{center}
\end{ejemplo}

%\pause

\textit{\underline{Solución}}:\\
\vspace*{15pt}
Un código cíclico debe ser lineal, y debe ser cerrado bajo la operación de desplazamiento cíclica. 

Por tanto, $C_1$ no es un código cíclico, debido a que \textbf{no es un código lineal} (100 + 010 = 110, el cual no está $C_1$, por ejemplo).\\
\vspace*{10pt}

$C_2$ es un \textbf{código cíclico}, debido a que se cumple ambas condiciones: por ejemplo, el desplazamiento cíclico de la palabra codificada 1010 es la palabra codificada 0101.

\end{frame} 

%%%%%%%%%%%%%%%%%%%%%%%%%%%%%%%%
\begin{frame}  \footnotesize
\begin{itemize} 

\item La expresión

\begin{center}
$a(x) = a_0 +a_1x+a_2x^2 + \ldots +a_dx^d$
\end{center}

en el cual $ a_0,a_1,a_2, \ldots, a_d$ están en $\mathbb{F}_2$, es denominado un \textbf{polinomio} con coeficientes en $\mathbb{F}_2$. Si $a_d \neq 0$ entonces el \textbf{grado} de a(x) es \textbf{d}. 

\item Los polinomios pueden ser sumados y multiplicados acorde con las reglas que conocemos de la algebra elemental y forma un \textbf{anillo}, $\mathbb{F}_2[x]$.

\item Tenemos una correspondencia entre polinomio $a_0 + a_1x +\ldots + a_{n-1} x^{n-1}$ en $\mathbb{F}_2[n]$, y la palabra $a_0 a_1\ldots a_{n-1}$ en $\mathbb{F}_2^n$.




\end{itemize} \end{frame} 

%%%%%%%%%%%%%%%%%%%%%%%%%%%%%%%%
\begin{frame}  \footnotesize  
\begin{itemize} 

\item El desplazamiento cíclico $\hat{a}$ de $\textbf{a}$ es representado por el polinomio $\hat{a}(x)$, donde:
\[
    \begin{array}{ll}
					  \hat{a}(x) &   = a_{n-1}+ a_0 x+ \ldots + a_{n-2} x^{n-1}  \\ 
					   &   = x(a_0 +a_1x+\ldots + a_{n-1}x^{n-1})-a_{n-1} (x^n -1)  \\ 
					   &   = x a(x) - a_{n-1}(x^n - 1)\footnote{El signo negativo es para mantener la claridad, pero debido a que los coeficientes pertenecen a $\mathbb{F}_2$ podemos reescribir como un signo más.} 
     \end{array}	    		
\]
 
\item $\hat{a}(x)$ es igual $xa(x)$, excepto para un múltiplo de $x^n -1$; en otro caso
\begin{center}
$\hat{a}(x)$ es igual a $xa(x)$ modulo $x^n - 1$.
\end{center}

\item Para \textbf{multiplicar} los polinomios, dado a(x) y b(x), los multiplicamos de la forma usual, y los reducimos modulo $x^n -1$ . $x^n -1$ equivalente al 0, y para la reducción es equivalente remplazar $x^n$ por $1$, $x^{n+1}$ por $x$, $x^{n+2}$ por $x^2$, y continuamos así. 

\item $V^n[x]$ el \textbf{anillo de polinomios} con coeficientes en $\mathbb{F}_2$, utilizando la multiplicación modulo $x^n-1$, por lo que tenemos una correspondencia biyectiva con $\mathbb{F}_2^n$, específicamente $a(x) \leftrightarrow a$. 

%\item Además, si a(x) y b(x) corresponde a \textbf{a} y \textbf{b}, entonces $a(x) + b(x)$ corresponde a \textbf{a + b} y $xa(x)$  corresponde a $\hat{a}$, el primer desplazamiento cíclico de $a$.

\end{itemize} \end{frame} 


%%%%%%%%%%%%%%%%%%%%%%%%%%%%%%%%
\begin{frame}  \small  

\begin{ejemplo}

Escribe los polinomios en $V^6[x]$ que corresponden a las palabras 110101 y 010110, y encontrar sus \textbf{productos} como elementos de $V^6[x]$.
\end{ejemplo}
%\pause
\textit{\underline{Solución}}:\\
\vspace*{5pt}

110101 es representado mediante $1 + x + x^3 + x^5$ y 010110 esta representado por $x+x^3 +x^4$. \\

Multiplicando y estableciendo $x_6=1, x_7=x, x_8=x^2, x_9=x^3$, y demás, obtenemos\\

\vspace*{5pt}
$(1+x+x^3 +x^5)(x+x^3 +x^4)$\\
$= (x+x^3 +x^4)+(x^2 +x^4 +x^5)+(x^4 +x^6 +x^7)+(x^6 +x^8 +x^9)$\\
$= x+x^2 +x^3 +x^4 +x^5 +x^7 +x^8 +x^9$\\
$= \cancel{x}+\cancel{x^2}+\cancel{x^3} +x^4 +x^5+\cancel{x}+ \cancel{x^2}+\cancel{x^3}$\\ 
$= x^4 +x^5.$\\


\end{frame} 

%%%%%%%%%%%%%%%%%%%%%%%%%%%%%%%%
\begin{frame}  \small  
\begin{itemize} 

%\item Mostraremos que un código cíclico en $\mathbb{F}_2^n$ corresponde a una especie particular de subconjunto de $V^n[x]$. 

\item Sea $R$ un \textbf{anillo} con la propiedad conmutativa en la multiplicación. 

\item Un subconjunto $S$ de $R$ es denominado como \textbf{ideal} si
\begin{enumerate}
\item  $a,b \in S \Rightarrow a+b \in S$ 
\item  $r\in R$ y $ a\in S \Rightarrow ra 	\in S$.
\end{enumerate}

\item  Un \textbf{ideal} S es cerrado bajo suma y bajo multiplicación por un elemento de $R$.

\begin{teorema}
Un \textbf{código binario} con palabras codificadas de longitud $n$ es cíclico sii corresponde a un \textbf{ideal} en $V^n[x]$.
\end{teorema}

\item \textcolor{red}{La construcción de códigos cíclicos de longitud n es equivalente a la construcción de ideales en $V^n[x]$.}

%\item Sea f(x) un polinomio en $V^n[x]$.  

%El conjunto de todos los múltiplos de f(x) en $V^n[x]$ es un \textit{ideal}, si a(x) y b(x) son múltiplos de f(x), por lo que a(x) + b(x) y  p(x) a(x) para cualquier p(x). 

\item El \textbf{ideal} generado mediante $f(x)$ es $\langle f (x)\rangle$.

\end{itemize} \end{frame} 

%%%%%%%%%%%%%%%%%%%%%%%%%%%%%%%%
%%%%%%%%%%%%%%%%%%%%%%%%%%%%%%%%
\begin{frame}  \footnotesize  

\begin{ejemplo}\label{ej:Ideal}

Construya el \textbf{ideal} generado mediante $f(x) = 1 + x^2$ en $V^3[x]$, y escriba el código correspondiente $C \subseteq \mathbb{F}_2^3$.

\end{ejemplo}

\textit{\underline{Solución}}:\\
Multiplicaremos f(x) por cada elemento p(x) de $V^3[x]$ y reduciendo modulo $x^3 - 1$, obtendremos la siguiente relación:
\begin{center}
 \begin{tabular}{ccl}
	    $\mathbb{F}_2^3$ & p(x) &  $p(x)f(x) mod (x^3-1)$  \\\hline
					 000 &  0 &     $\mathbf{0}$\\ 
					 100 &  1 &     $\mathbf{1+x^2}$\\ 
					 010 &  x &     $x+x^3=\mathbf{1+x}$\\ 
					 110 &  $x+1$ &   $x+x^3+1+x^2=\mathbf {x+x^2} $ \\ 
					 001 &  $x^2$ &    $x^2 + x^4 = \mathbf {x^2 +x} $\\ 
					 101 &  $x^2+1$ &  $1+x^2+x^4+x^2 =\mathbf {1+x} $  \\ 
					 011 &  $x^2+x$ &   $ x^2+x+x^4+x^3=\mathbf {x^2+1} $\\ 
					 111 &  $x^2+x+1$ &   $x^2+x+1+x^4+x^3+x^2=\mathbf {0}$  \\ 

     \end{tabular}	    		
\end{center}

Por tanto el \textbf{ideal} es $\langle 1 + x^2 \rangle$ que tiene solamente cuatro elementos:

\begin{center}
$0$, $(x+1)$, $(x^2+1)$, $(x^2+x)$
\end{center}

El \textbf{código correspondiente} en $\mathbb{F}_2^3$ es $C = \{000, 011, 101, 110\}.$

\end{frame} 

%%%%%%%%%%%%%%%%%%%%%%%%%%%%%%%%
%%%%%%%%%%%%%%%%%%%%%%%%%%%%%%%%
%\begin{frame}  \small \begin{ejercicio} \end{ejercicio} \pause \textit{\underline{Solución}}:\\ \end{frame} 
%%%%%%%%%%%%%%%%%%%%%%%%%%%%%%%%
%%%%%%%%%%%%%%%%%%%%%%%%%%%%%%%%
%EJERCICIO %%%%%%%%%%%%%%%%%%%%%%%%%%%%%%%
%%%%%%%%%%%%%%%%%%%%%%%%%%%%%%%%

%%%%%%%%%%%%%%%%%%%%%%%%%%%%%%%%
%%%%%%%%%%%%%%%%%%%%%%%%%%%%%%%%
\begin{frame}  \small \begin{ejercicio} 

¿Cual de los siguientes códigos son cíclicos?

$C_1 = \{0000, 1100, 0110, 0011, 1001\},$\\
$C_2 = \{0000, 1100, 0110, 0011, 1001, 1010, 0101, 1111\}.$\\

\end{ejercicio} %\pause 
\textit{\underline{Solución}}:\\ 
 \vspace*{5pt}
$C_1$:1100+0110 = 1010 No es lineal.\\
\vspace{0.3cm}

$C_2$: es cíclico \\
\vspace{0.3cm}
0110+0011= 0101,  0110+1001=1111, 0110+1010=1100, 0110+0101=0011.


\end{frame} 

%%%%%%%%%%%%%%%%%%%%%%%%%%%%%%%%
%%%%%%%%%%%%%%%%%%%%%%%%%%%%%%%%
\begin{frame}  \small \begin{ejercicio} 

Escribe las palabras codificadas del código cíclico correspondiente al ideal $\langle 1+x+x^2\rangle$, en $V^3[x]$ y encuentra la matriz de verificación para este código.

\end{ejercicio} 

 \textit{\underline{Solución}}:\\ 
 \vspace*{5pt}
Multiplicaremos f(x) por cada elemento p(x) de $V^3[x]$ y reduciendo modulo $x^3 - 1$, obtendremos la siguiente relación:
\begin{center}

\[
    \begin{array}{ccl}
					 \mathbb{F}_2^3 & p(x) &  p(x)f(x) mod (x^3-1)  \\ \hline
					 000 &  0 &     \mathbf{0}\\ 
					 100 &  1 &     \mathbf{1+x+x^2}\\ 
					 010 &  x &     x+x^2+x^3=\mathbf{1+x+x^2}\\ 
					 110 &  x+1 &   x+x^2+x^3+1+x+x^2 =\mathbf {0}  \\ 
					 001 &  x^2 &    x^2  +x^3+ x^4 = \mathbf {1+x+x^2} \\ 
					 101 &  x^2+1 & x^2+x^3+x^4+1+x+x^2 =\mathbf {0}   \\
					 011 &  x^2+x &    x^2+x^3+x^4+x+x^2+x^3=\mathbf {0} \\  
					 111 &  x^2+x+1 &   x^2+x^3+x^4+x+x^2+x^3+1+x+x^2=\mathbf {1+x+x^2}  

     \end{array}	    		
\]
\end{center}

\end{frame} 

%%%%%%%%%%%%%%%%%%%%%%%%%%%%%%%%
\begin{frame}  \small

Por tanto el ideal es $\langle 1 +x+ x^2 \rangle$ que tiene solamente 2 elementos:

\begin{center}
$0$,$(1+x+x^2)$
\end{center}

El código correspondiente en $\mathbb{F}_2^3$ es $C = \{000, 111 \}.$\\
\vspace*{5pt}

Los parámetros son n=3; $n=2^m -1$; m=2. $k =n-m =3-2=1$\\
\vspace*{5pt}

La matriz es: 

\[
    H=   \left( \begin{array}{ccc}
0&	1	& 1	\\		
1&	0	& 1	\\
           			\end{array}
    		\right) 	    		
\]



\end{frame} 

 


%9.7. Does the result in the previous exercise hold when 5 is replaced by an arbitrary integer n ! 2?


%%%%%%%%%%%%%%%%%%%%%%%%%%%%%%%%
%%%%%%%%%%%%%%%%%%%%%%%%%%%%%%%%
%SECCION %%%%%%%%%%%%%%%%%%%%%%%%%%%%%%%
%%%%%%%%%%%%%%%%%%%%%%%%%%%%%%%%
%%%%%%%%%%%%%%%%%%%%%% 


\begin{frame}   \footnotesize
{
\setbeamercolor{block title}{bg=blue!60!black, fg=white}
\setbeamercolor{block body}{bg=blue!10, fg=black}
\begin{block}{Sección 2.1}
Propiedades de los códigos cíclicos\end{block}
}
\end{frame} 


\subsection{Propiedades de los códigos cíclicos}

%%%%%%%%%%%%%%%%%%%%%%%%%%%%%%%%
\begin{frame}  \frametitle{Propiedades de los códigos cíclicos}\small  

El próximo teorema muestra que todos los códigos cíclicos corresponden a un ideal generado por un polinomio.

\begin{teorema}\label{te:ideal}

Sea $C \neq \{0\}$ un código cíclico (ideal) en $V^n[x]$. Entonces existe un polinomio g(x) en C tal que $C=\langle g(x)\rangle$.

\end{teorema}


El ejemplo \ref{ej:Ideal} muestra el ideal en $V^3[x]$ generado por $1+x^2$. 


\begin{itemize} 

\footnotesize

\item El polinomio de menor grado, que no es cero, es 1+x, y el teorema \ref{te:ideal} nos dice que este polinomio es \textbf{también un generador para el ideal}. 

\item Un código cíclico C tendrá numerosos generadores, pero solamente uno de ellos tendrá el menor grado en C y será el \textbf{generador canónico} de C.


\end{itemize} 


\end{frame} 


%%%%%%%%%%%%%%%%%%%%%%%%%%%%%%%%
\begin{frame}  \small  

\begin{teorema}\label{te:GeneradorCanonico}
El generador canónico g(x) de un código cíclico en $V^n[x]$ es un divisor de $x^n -1$ en $\mathbb{F}_2[x]$.
\end{teorema}

El generador canónico g(x) determina la \textbf{dimensión} de un código cíclico, y la \textbf{matriz de verificación} asociada.

\footnotesize
\begin{itemize} 

\item   Ya que  g(x) es un divisor $x^n-1$, tendremos que $g(x)h(x) = x^n-1$ en $\mathbb{F}_2[x]$.


\item Sea

\begin{center}
$g(x)=g_0 +g_1x+\ldots +g_{n-k} x^{n-k}$,	$h(x)=h_0 +h_1x+\ldots+h_kx^k$,
\end{center}

donde los coeficiente s $g_0$, $h_0$, $g_{n-k}$, $h_k$ deben ser todos 1, ya que el producto de los polinomios es $x^n - 1$. 

\item Sea \textbf{g} la palabra en $\mathbb{F}_2^n$

\begin{center}
$g = g_0 g_1 \ldots g_{n-k}0 0\ldots 0,$
\end{center}

\item Sea $\mathbf{h^*}$ la palabra aquella cuales los primeros $k+1$ bits son los coeficientes de h(x) a la inversa, seguidos de n - k - 1 ceros:

\begin{center}
$h^* = h_k h_{k-1} \ldots h_00 0\ldots 0,$
\end{center}

\end{itemize} \end{frame} 


%%%%%%%%%%%%%%%%%%%%%%%%%%%%%%%%
\begin{frame}  \footnotesize  

Sea \textbf{H} la matriz (n-k) x n cuyas filas estan formadas primero por $h^*$ y las siguientes filas por n-k-1 desplazamientos cíclicos de $h^*$:


\[
    H =   \left( \begin{array}{ccccccc}
     				  h_k &  h_{k-1} & \ldots   &   h_0  &   0  &	0    & 0\\ 
					  0     &   h_k    & \ldots		&   h_1  & h_0 &	0    & 0\\ 
					  0     &   0        & \ldots     &   h_2  & h_1 &	h_0 & 0\\
  					 .      &   .        & \ldots     &   .  & . &	. & .\\
					 0     &   0           & \ldots  &   h_k  & h_{k-1} &	\ldots& h_0
           			\end{array}
    		\right) 	    		
\]

\begin{lema}
Sea g(i) la fila vector  $0 0\ldots 0 g_0g_1 \ldots g_{n-k} 0 0\ldots 0$ donde hay $i$ ceros al comienzo y k-1-i ceros al final, o sea, el desplazamiento cíclico de $g$ correspondiente a $x^i g(x)$ en $V^n[x]$. Entonces $Hg'_{(i)}= 0'$ $(0 \leq i \leq n - 1).$
\end{lema}


\begin{teorema}\label{te:DimensionC}
La matriz \textbf{H} es una \textbf{matriz de verificación} para el código cíclico $C=\langle g(x) \rangle$, y la \textbf{dimensión} de C es \textbf{k}.
\end{teorema}
\vspace*{5pt}

\textcolor{red}{Para encontrar todos los códigos cíclicos con longitud de palabra n, es suficiente encontrar los \textbf{factores irreducibles} de $x^n-1$ en $\mathbb{F}_2[x]$}. 

\end{frame} 
%%%%%%%%%%%%%%%%%%%%%%%%%%%%%%%%

%%%%%%%%%%%%%%%%%%%%%%%%%%%%%%%%
\begin{frame}  \small  
\begin{ejemplo}

Dado que hay tres factores irreducible de $x^7-1$ en $\mathbb{F}_2[x]$:

\begin{center}
$x^7 -1$=$(1+x)$ $(1+x+x^3)$ $(1+x^2 +x^3)$,
\end{center}

¿Cuales son las posibilidades para un código cíclico con longitud de palabra de 7?\\
\end{ejemplo}


\textit{\underline{Solución}}:\\
\vspace*{5pt}
Combinando los tres factores irreducibles en todas las posibilidades podemos obtener $2^3=8$ divisores de $x^7-1$ en $\mathbb{F}_2[x]$.\\
\vspace*{5pt}
Los divisores son además de $1$ y $x^7 - 1$,

\begin{center}
$(1+x)$,	$(1+x+x^3)$,	$(1+x^2 +x^3)$,	$(1+x)(1+x+x^3)$,\\
$(1+x)$ $(1+x^2 +x^3)$,	 $(1+x+x^3)$ $(1+x^2 +x^3)$.

\end{center}

Cada uno de los divisores generan un código cíclico, y  (por teorema \ref{te:ideal} y \ref{te:GeneradorCanonico}) son los únicos códigos cíclicos de longitud 7. \\
\vspace*{5pt}
$\langle 1\rangle$ es el código en el cual cada palabra es una palabra codificada, y $\langle x^7 -1\rangle = \langle 0 \rangle$ es el código en el cual la única palabra codificada es 0. \\
\vspace*{5pt}
\end{frame} 

%%%%%%%%%%%%%%%%%%%%%%%%%%%%%%%%
\begin{frame}  \small  

Los otros códigos son más interesantes. \\
\vspace*{5pt}


Por ejemplo, sea C el código con un generador canónico $g(x) = 1 + x + x^3$, por lo que 

\begin{center}
$h(x)=(1+x)(1+x^2 +x^3)=1+x+x^2 +x^4$
\end{center}

y $h^* = 1011100$. Del teorema \ref{te:DimensionC} obtenemos que la dimensión de C es 4, y la matriz de verificación para C es:


\[
    H =   \left( \begin{array}{ccccccc}
     				  1 &  0  & 1  & 1 & 1  & 0  & 0\\ 
					  0 &  1  & 0	&  1 & 1 &	 1  & 0\\ 
					  0 &  0  & 1  &  0 & 1 & 1  & 1
           			\end{array}
    		\right) 	    		
\]

Las sietes columnas de esta matriz son distintas y no son iguales a cero, por lo que el código C es equivalente al código Hamming en $\mathbb{F}_2^7$.

\end{frame} 
%%%%%%%%%%%%%%%%%%%%%%%%%%%%%%%%
\begin{frame}  \tiny  
$C=\langle g(x) =1 +x +x^3 \rangle$, n=7, entonces k=n-grado g; k=7-3=4\\
\vspace{0.3cm}
Los mensajes posibles son todos los polinomios de grado menor que $4$:
\vspace{0.3cm}

$m(x)\in\{\,0,\ 1,\ x,\ x^{2},\ x^{3},\ 
1+x,\ 1+x^{2},\ 1+x^{3},\ 
x+x^{2},\ x+x^{3},\ x^{2}+x^{3},\
1+x+x^{2},\ 1+x+x^{3},\ 1+x^{2}+x^{3},\ 
x+x^{2}+x^{3},\ 1+x+x^{2}+x^{3}\,\}.
$\\

\vspace{0.3cm}

Cada palabra del código es
\[
c(x)=m(x)\,g(x)\ \bmod (x^{7}-1).
\]

A continuación se muestran las $16$ palabras del código $C$ en forma de
vectores de longitud $7$:

\[
\begin{array}{c|c}
m(x) & c(x) \\ \hline
0 & 0000000 \\
1 & 1101000 \\
x & 0110100 \\
x^{2} & 0011010 \\
x^{3} & 0001101 \\
1+x & 1011100 \\
1+x^{2} & 1110010 \\
1+x^{3} & 1100101 \\
x+x^{2} & 0101110 \\
x+x^{3} & 0111001 \\
x^{2}+x^{3} & 0010111 \\
1+x+x^{2} & 1000110 \\
1+x+x^{3} & 1010011 \\
1+x^{2}+x^{3} & 1111111 \\
x+x^{2}+x^{3} & 0100011 \\
1+x+x^{2}+x^{3} & 1001011 \\
\end{array}
\]


\end{frame} 


%%%%%%%%%%%%%%%%%%%%%%%%%%%%%%%%
%\begin{frame}  \small  
%\begin{itemize} \item  \item \item \item 
%\end{itemize} \end{frame} 


%%%%%%%%%%%%%%%%%%%%%%%%%%%%%%%%
%\begin{frame}  \small  
%\end{frame} 
%%%%%%%%%%%%%%%%%%%%%%%%%%%%%%%%

%%%%%%%%%%%%%%%%%%%%%%%%%%%%%%%%
%%%%%%%%%%%%%%%%%%%%%%%%%%%%%%%%
%\begin{frame}  \small \begin{ejercicio} \end{ejercicio} \pause \textit{\underline{Solución}}:\\ \end{frame} 
%%%%%%%%%%%%%%%%%%%%%%%%%%%%%%%%
%%%%%%%%%%%%%%%%%%%%%%%%%%%%%%%%
%EJERCICIO %%%%%%%%%%%%%%%%%%%%%%%%%%%%%%%
%%%%%%%%%%%%%%%%%%%%%%%%%%%%%%%%

%%%%%%%%%%%%%%%%%%%%%%%%%%%%%%%%
%%%%%%%%%%%%%%%%%%%%%%%%%%%%%%%%
\begin{frame}  \small \begin{ejercicio} 

Escribe los factores $x^3-1$ en $\mathbb{F}_2[x]$ y determina todos los códigos cíclicos de longitud 3.

\end{ejercicio} %\pause 
\textit{\underline{Solución}}:\\ 

Sabemos que los códigos cíclicos corresponden a los ideales en $V^3[x]$ y sabemos que $X^3-1=(X+1) (X^2+X+1)$ en $\mathbb{F}_2[x]$. Por tanto tenemos 4 códigos cíclicos:

\begin{itemize}
\item $g=1$: que obtenemos $C=\mathbb{F}_2^3$
\item $g=X+1$: donde $C=\{000,011,101,110\}$ que es el código de verificación de paridad (el conjunto de todas las palabras de peso par)
\item $g=X^2+X+1$: que es $C=\{000,111\}$ el código de repetición triple.
\item $g=X^3-1$: donde $C=\{000\}$
\end{itemize}


\end{frame} 



%%%%%%%%%%%%%%%%%%%%%%%%%%%%%%%%
\begin{frame}  \small \begin{ejercicio} 

 Completa la clasificación de los códigos cíclicos de longitud 7, mediante las lineas del Ejemplo 5.


\end{ejercicio} %\pause 
\textit{\underline{Solución}}:\\ 

\begin{itemize}
\item En el caso de g(x)= $1+x$ tenemos h(x)= $(1+x^2+x^3)$ $(1+x+x^3)$= $1+x+x^2+x^3+x^4+x^5+x^6$

y $h^* =1111111$. Sabemos que la dimensión de C (k) =1, por tanto la matriz de verificación.

\[
    H =   \left( \begin{array}{ccccccc}
     				  1 &  1  & 1  & 1 & 1  & 1  & 1
           			\end{array}
    		\right) 	    		
\]

\vspace*{5pt}
$C=\{$ 1100000, 0110000, 0011000, 0001100, 0000110, 0000011, 1000001, 0000000$\}$\\
\vspace*{5pt}

\end{itemize}

\end{frame} 

%%%%%%%%%%%%%%%%%%%%%%%%%%%%%%%%
\begin{frame}  \small 


\begin{itemize}
\item En el caso de g(x)= $1+x^2+x^3$ tenemos h(x)= $(1+x)$ $(1+x+x^3)$= $1+x^2+x^3+x^4$
y $h^* =1110100$. Sabemos que la dimensión de C (k) =4, por tanto la matriz de verificación.

\[
    H =   \left( \begin{array}{ccccccc}
     				  1 &  1  & 1  & 0 & 1  & 0  & 0\\ 
					  0 &  1  & 1	&  1 & 0 &	 1  & 0\\ 
					  0 &  0  & 1  &  1 & 1 & 0  & 1
           			\end{array}
    		\right) 	    		
\]

El código resultante es un código Hamming.

$C=\{$ 1011000, 0101100, 0010110, 0001011, 1000101, 1100010, 0110001,0000000$\}$\\
\vspace*{5pt}

\end{itemize}

\end{frame} 
%%%%%%%%%%%%%%%%%%%%%%%%%%%%%%%%
\begin{frame}  \small

\begin{itemize}
\item En el caso de g(x)= $(1+x)$ $(1+x^2+x^3)$=$1+x+x^2+x^4$ tenemos h(x)= $(1+x+x^3)$
y $h^* =1011000$. Sabemos que la dimensión de C (k) =3, por tanto la matriz de verificación.

\[
    H =   \left( \begin{array}{ccccccc}
     				  1 &  0  & 1  & 1 & 0  & 0  & 0\\ 
					  0 &  1  & 0	&  1 & 1 &	 0  & 0\\ 
					  0 &  0  & 1  &  0 & 1 & 1  & 0\\
					  0 &  0  & 0  &  1 & 0 & 1  & 1
           			\end{array}
    		\right) 	    		
\]

$C=\{$ 1110100, 0111010, 0011101, 1011100, 0101110, 0010111, 1010011, 1101001, 0000000$\}$
\end{itemize}

\end{frame} 

%%%%%%%%%%%%%%%%%%%%%%%%%%%%%%%%
%%%%%%%%%%%%%%%%%%%%%%%%%%%%%%%%

\begin{frame}  \small

\begin{itemize}
\item En el caso de g(x)= $(1+x)(1+x+x^3)$=$1+x^2+x^3+x^4$ tenemos h(x)= $(1+x^2+x^3)$
y $h^* =1101000$. Sabemos que la dimensión de C (k) =3, por tanto la matriz de verificación.

\[
    H =   \left( \begin{array}{ccccccc}
     				  1 &  1  & 0  & 1 & 0  & 0  & 0\\ 
					  0 &  1  & 1	&  0 & 1 &	 0  & 0\\ 
					  0 &  0  & 1  &  1 & 0 & 1  & 0\\
					  0 &  0  & 0  &  1 & 1 & 0  & 1
           			\end{array}
    		\right) 	    		
\]

\vspace*{5pt}
$C=\{$ 1011100, 0101110, 0010111, 1001011, 1100101, 1100010, 0110001,000000$\}$
\vspace*{5pt}

\end{itemize}

\end{frame} 


%%%%%%%%%%%%%%%%%%%%%%%%%%%%%%%%
\begin{frame}  \small 


\begin{itemize}
\item En el caso de g(x)= $(1+x^2+x^3) (1+x+x^3)=1+x+x^2+x^3+x^4+x^5+x^6$ tenemos h(x)= $(1+x)$
y $h^* =1100000$. Sabemos que la dimensión de C (k) =1, por tanto la matriz de verificación.

\[
    H =   \left( \begin{array}{ccccccc}
     				  1 &  1  & 0  & 0 & 0  & 0  & 0\\ 
					  0 &  1  & 1	&  0 & 0 &	 0  & 0\\ 
					  0 &  0  & 1  &  1 & 0 & 0  & 0\\
					  0 &  0  & 0  &  1 & 1 & 0  & 0\\
					  0 &  0  & 0  &  0 & 1 & 1  & 0\\
					  0 &  0  & 0  &  0 & 0 & 1  & 1
           			\end{array}
    		\right) 	    		
\]


\vspace*{5pt}
$C=\{$ 1111111,000000$\}$


\end{itemize}

\end{frame} 
%9.10. Show that when n is odd, x − 1 occurs exactly once as a factor of xn − 1 in F2[x]. (It is worth remembering that x − 1 and 1+x are the same in F2[x].)


% 9.12. Write down the equations for the coefficients of g(x) and h(x) that result from the fact that g(x)h(x) = xn−1, and hence prove Lemma 9.12


%%%%%%%%%%%%%%%%%%%%%%%%%%%%%%%%
%%%%%%%%%%%%%%%%%%%%%%%%%%%%%%%%
\begin{frame}  \small \begin{ejercicio} 

Considere el generador $g(X)= X^3 + X +1$ en $V^7[x]$ (sabiendo que g divide $X^7-1$). Codifique el mensaje m=1010 y decodifique el mensaje z= 1001011.


\end{ejercicio} %\pause 
\textit{\underline{Solución}}:\\ 

\begin{description}

\item [Codificación]


Transformamos el mensaje 1010 en un polinomio de la forma m(X)=$X^2+1$ y lo multiplicamos por g(X) y obtenemos c(X)=$m(X)g(X)$ = $X^5+2X^3+X^2+X+1$ que nos da la palabra codificada 1110010.

Hubiéramos obtenido el mismo mensaje si en vez de utilizar el polinomio utilizáramos la matriz generadora.

\[
    G =   \left( \begin{array}{ccccccc}
     				  1 &  1  & 0  & 1 & 0  & 0  & 0\\ 
					  0 &  1  & 1	&  0 & 1 &	 0  & 0\\ 
					  0 &  0  & 1  &  1 & 0 & 1  & 0\\
					  0 &  0  & 0  &  1 & 1 & 0  & 1
           			\end{array}
    		\right) 	    		
\]

\end{description}

\end{frame} 


%%%%%%%%%%%%%%%%%%%%%%%%%%%%%%%%
\begin{frame}  \footnotesize 

 \textit{\underline{cont Solución}}:\\ 

\begin{description}

\item [Decodificación]

Sabemos que la matriz de verificación es la siguiente.
\[
    H =   \left( \begin{array}{ccccccc}
     				  0 &  0  & 1  & 0 & 1  & 1  & 1\\ 
					  0 &  1  & 0	&  1 & 1 &	 1  & 0\\ 
					  1 &  0  & 1  &  1 & 1 & 0  & 0
           			\end{array}
    		\right) 	    		
\]

Si multiplicamos Hz = (1001011) H =(000) por lo que pertenece a la codificación.

Sabemos que $z(X) = c(X) g(X)$, por lo que solamente tenemos que buscar el polinomio resto de la división por el polinomio generador. 

$1+X^3+X^5+X^6 = c(X) (X^3+X+1)$ por tanto \\
\vspace{0.3cm}
\tiny
$\textcolor{orange}{\mathbf{X^3}} * (X^3+X+1) = \textcolor{red}{X^6 +X^4 + X^3} $\\
\hspace{0.3cm}$ 1+X^3+X^5+X^6 +\textcolor{red}{X^6 +X^4 + X^3} = \textcolor{blue}{X^5 + X^4 +1}$

$\textcolor{orange}{\mathbf{X^2}} * (X^3+X+1) = \textcolor{olive}{X^5+X^3+X^2} $\\
\hspace{0.3cm}$\textcolor{blue}{X^5 + X^4 +1} + \textcolor{olive}{X^5+X^3+X^2} = \textcolor{brown}{X^4 + X^3+X^2 +1}   $

$\textcolor{orange}{\mathbf{X}} * (X^3+X+1) =\textcolor{magenta}{ X^4+X^2+X }$\\
\hspace{0.3cm}$\textcolor{brown}{X^4 + X^3+X^2 +1} +\textcolor{magenta}{ X^4+X^2+X } =   X^3 + X + 1$\\
 
$\textcolor{orange}{\mathbf{1}} * (X^3+X+1) = \textcolor{pink}{X^3 + X + 1}$\\
\hspace{0.3cm} $ X^3 + X + 1 + \textcolor{pink}{X^3 + X + 1}= 0$\\
\vspace{0.3cm}
\footnotesize
$c(X) = \textcolor{orange}{1+X+ X^2+X^3}$ y c=(1111).


\end{description}

\end{frame}
%%%%%%%%%%%%%%%%%%%%%%%%%%%%%%%%
%%%%%%%%%%%%%%%%%%%%%%%%%%%%%%%%
%%%%%%%%%%%%%%%%%%%%%%%%%%%%%%%%

\begin{frame}   \footnotesize
{
\setbeamercolor{block title}{bg=blue!60!black, fg=white}
\setbeamercolor{block body}{bg=blue!10, fg=black}
\begin{block}{Sección 2.2}
Corrección Errores en los códigos cíclicos\end{block}
}
\end{frame} 


\subsection{Corrección Errores en los códigos cíclicos}

%%%%%%%%%%%%%%%%%%%%%%%%%%%%%%%%
\begin{frame}  \frametitle{Corrección Errores en los códigos cíclicos}\small  

%Sea C un código cíclico corrector de error \textcolor{red}{$t$} y el mensaje recibido es \textcolor{red}{w(X)}.


Sea C un código cíclico de longitud n corrector de error \textcolor{red}{$t$}, con polinomio generador g(X). \\
\vspace{0.3cm}
Recibimos una palabra w(X) que puede tener errores, w(X)=c(X)+e(X) donde:
\begin{itemize}
    \item c(X) es el código correcto (múltiplo de g(X))
    \item e(X) es el polinomio de error.
\end{itemize}




\begin{enumerate} 

\item  Calcular el polinomio \textbf{síndrome} $s(X) \equiv w(X)\, mod\, g(X)$\\
\vspace{0.3cm}
Como c(X) es múltiplo de g(X), su residuo es 0.\\

Entonces $s(X)= e(X)\, mod\, g(X)$


\item Para $i=0, \ldots, n-1$ calcula los \textbf{desplazamientos cíclicos} $s_i(X) \equiv X^i s(X)\, mod\, g(X)$ hasta encontrar un \textbf{síndrome} $s_j$ tal que $w(s_j)\leq t$. \\
\vspace*{5pt}

\textcolor{red}{La palabra error de menor peso es:} 
\begin{center}
$e(X) =X^{n-j} s_j(X)\, mod\, (X^n-1)$
\end{center}


\end{enumerate} 


\end{frame} 

%%%%%%%%%%%%%%%%%%%%%%%%%%%%%%%%
\begin{frame}  \footnotesize  
\begin{ejemplo}

Sea $g(X) = X^3 + X+1$ y w=(1011011);


\end{ejemplo}
 \textit{\underline{Solución}}:\\ 
 \vspace{0.2cm}
$ g(X) = X^3 + X + 1, \qquad 
w=(1011011) \rightarrow w(X)= 1+X^2+X^3+X^5+X^6$
\vspace{0.2cm}\\
Realizamos la división en $\mathbb{F}_2[X]$ (sumas = XOR):
\vspace{0.2cm}\\

\textbf{Paso 1:} $X^6 / X^3 = \mathbf{X^3}  \Rightarrow \text{ restamos } X^3 g(X) = X^6 + X^4 + X^3$, \\
$R_1 = (X^6 + X^5 + X^3 + X^2 + 1) \oplus (X^6 + X^4 + X^3) = X^5 + X^4 + X^2 + 1.$ \\[3pt]

\textbf{Paso 2:}  $X^5 / X^3 = \mathbf{X^2}$ 
$\Rightarrow \text{ restamos } X^2 g(X) = X^5 + X^3 + X^2$, \\[4pt]
$R_2 = (X^5 + X^4 + X^2 + 1) \oplus (X^5 + X^3 + X^2) = X^4 + X^3 + 1. $\\[3pt]

\textbf{Paso 3:}  $X^4 / X^3 = \mathbf{X} $
$\Rightarrow \text{ restamos } X g(X) = X^4 + X^2 + X,$ \\[4pt]
$R_3 = (X^4 + X^3 + 1) \oplus (X^4 + X^2 + X) = X^3 + X^2 + X + 1. $\\[3pt]

\textbf{Paso 4:} $X^3 / X^3 = \mathbf{1} $
$\Rightarrow \text{ restamos } g(X) = X^3 + X + 1,$ \\[4pt]
$R_4 = (X^3 + X^2 + X + 1) \oplus (X^3 + X + 1) = \boxed{X^2}.$\\

\vspace{0.2cm}

Como $\deg(R_4) = 2 < 3$, la división termina.
\[
\boxed{w'(X) = (X^3 + X^2 + X + 1)\,g(X) + X^2}
\]

\end{frame}
%%%%%%%%%%%%%%%%%%%%%%%%%%%%%%%%
%%%%%%%%%%%%%%%%%%%%%%%%%%%%%%%%

\begin{frame}  \footnotesize
Por tanto:
\[
q(X) = X^3 + X^2 + X + 1, 
\qquad
r(X) = X^2.
\]

\vspace{0.5cm}

Asumiendo que solamente ha ocurrido un error:
\vspace{0.2cm}
w($X^2)= 1 <= t = 1$\\
\vspace{0.2cm}
la palabra corregida es (10\textbf{0}1011) de la  palabra recibida w=(10\textbf{1}1011) .

\end{frame} 

%%%%%%%%%%%%%%%%%%%%%%%%%%%%%%%%
\begin{frame}  \footnotesize
\begin{ejemplo}

Sea $g(X) = X^8 + X^7+ X^6+ x^4+ 1$ un generador de tipo [15,7,5], o lo que es lo mismo un código corrector 2. Si recibimos w=(110011101100010); corrige el posible error.


\end{ejemplo}
 %\pause 
 \textit{\underline{Solución}}:\\ 

 $ g(X) = X^8 + X^7+ X^6+ x^4+ 1, \qquad 
w=(110011101100010) \rightarrow w(X)= 1+X+X^4+X^5+X^6+X^8+X^9+X^{13}$
    
\vspace{0.3cm}
\textbf{Paso 1: } $ X^{13} / X^8 =\mathbf{X^5} \;\Rightarrow\; X^{5}g = X^{13}+X^{12}+X^{11}+X^{9}+X^{5}$
\\
$R_1 = X^{13}+X^{9}+X^{8}+X^{6}+X^{5}+X^{4}+X+1  \oplus X^{5}g = X^{12}+X^{11}+X^{8}+X^{6}+X^{4}+X+1$ \\[6pt]


\textbf{Paso 2: } $X^{12}/X^{8}= \mathbf{X^4} \;\Rightarrow\;
X^{4}g = X^{12}+X^{11}+X^{10}+X^{8}+X^{4}$ \\
$R_2 = R_1 \oplus X^{4}g
     = X^{10}+X^{6}+X+1$ \\[6pt]


\textbf{Paso 3: } $X^{10} / X^8=\mathbf{X^2} \;\Rightarrow\;
X^{2}g = X^{10}+X^{9}+X^{8}+X^{6}+X^{2}$ \\
$R_3 = R_2 \oplus X^{2}g
     = X^{9}+X^{8}+X^{2}+X+1$ \\[6pt]
     
\textbf{Paso 4: }  $X^9 / X^8= \mathbf{X}
\;\Rightarrow\;
Xg = X^{9}+X^{8}+X^{7}+X^{5}+X $\\
$R_4 = R_3 \oplus Xg
     = \boxed{\,X^{7}+X^{5}+X^{2}+1\,}$

\vspace{0.3cm}

Entonces $w(X)= (X+X^2+X^4+X^5) g(X)+ (1+X^2+X^5+X^7)$


\end{frame} 
%%%%%%%%%%%%%%%%%%%%%%%%%%%%%%%%
\begin{frame}  \footnotesize  

$\mathbf{s_0(X)} =  (1+X^2+X^5+X^7) = \boxed{10100101}$ \textcolor{blue}{$w(s_0) = 4.$}
sigo calculando síndromes hasta $w(s_j)\leq t$:\\
\vspace{0.3cm}
$\mathbf{s_1(X)}= X s_0 (X) \; mod \; g$
$=X(1+X^2+X^5+X^7)= X+X^3+X^6+\textcolor{red}{X^8} =  X+X^3+X^6+\textcolor{red}{X^7+X^6+X^4+1}$ = \boxed{1+X+X^3+X^4+X^7} = \boxed{11011001} \textcolor{blue}{$w(s_1) = 5.$}

\vspace{0.3cm}
$\mathbf{s_2(X)}= X s_1 (X) \; mod \; g$
$=X(1+X+X^3+X^4+X^7)= X+X^2+X^4+X^5+\textcolor{red}{X^8} =  X+X^2+X^4+X^5+\textcolor{red}{X^7+X^6+X^4+1}$ = \boxed{1+X+X^2+X^5+X^6+X^7} = \boxed{11100111}\textcolor{blue}{$w(s_2) = 6.$}

\vspace{0.3cm}
$\mathbf{s_3(X)}= X s_2 (X) \; mod \; g$
$=X(1+X+X^2+X^5+X^6+X^7)= X+X^2+X^3+X^6+X^7+\textcolor{red}{X^8} =  X+X^2+X^3+X^6+X^7+\textcolor{red}{X^7+X^6+X^4+1}$ = \boxed{1+X+X^2+X^3+X^4} = \boxed{11111000} \textcolor{blue}{$w(s_3) = 5.$}

\vspace{0.3cm}
$\mathbf{s_4(X)}= X s_3 (X) \; mod \; g$
$=X(1+X+X^2+X^3+X^4)= \boxed{X+X^2+X^3+X^4+X^5} =  \boxed{01111100}$ \textcolor{blue}{$w(s_4) = 5.$}

\vspace{0.3cm}
$\mathbf{s_5(X)}= X s_4 (X) \; mod \; g$
$=X(X+X^2+X^3+X^4+X^5)= \boxed{X^2+X^3+X^4+X^5+X^6} =  \boxed{001111110}$ \textcolor{blue}{$w(s_5) = 5.$}

\end{frame} 
%%%%%%%%%%%%%%%%%%%%%%%%%%%%%%%%
\begin{frame}  \footnotesize  

\vspace{0.3cm}
$\mathbf{s_6(X)}= X s_5 (X) \; mod \; g$
$=X(X^2+X^3+X^4+X^5+X^6)=  \boxed{X^3+X^4+X^5+X^6+X^7} =  \boxed{000111111}$ \textcolor{blue}{$w(s_6) = 5.$}

\vspace{0.3cm}
$\mathbf{s_7(X)}= X s_6 (X) \; mod \; g$
$=X(X^3+X^4+X^5+X^6+X^7)=  X^4+X^5+X^6+X^7+\textcolor{red}{X^8} =  X^4+X^5+X^6+X^7+\textcolor{red}{X^7+X^6+X^4+1}$ = \boxed{1+X^5} = \boxed{10000100} \textcolor{blue}{$w(s_7) = 2.$}
\[
     \begin{array}{c|c}
     			\tiny	  s_0 & 10100101\\ 
     				  s_1 & 11011001\\ 
     				  s_2 & 11100111\\ 
     				  s_3 & 11111000\\ 
     				  s_4 & 01111100\\ 
     				  s_5 & 00111110\\ 
     				  s_6 & 00011111\\ 
     				  s_7 & 10000100 
           			\end{array}
     \]
\begin{center}
$\mathbf{e(X)} =X^{n-j} s_j(X)\, mod\, (X^n-1)$, $e(X) =X^{15-\mathbf{7}} s_\mathbf{7}(X)\, mod\, (X^{15}-1)$\\
\vspace{0.2cm}
$\mathbf{e(X)} =X^8 (1+X^5) \, mod\, (X^{15}-1)= X^8 + X^{13} \rightarrow \boxed{e = 000000001000010}$

\vspace{0.2cm}

$\mathbf{r(X)}=w(X)+e(X) = (110011101100010) + (000000001000010) = $

\vspace{0.1cm}
$\boxed{(110011100100000)}$
\end{center}


\end{frame} 


%%%%%%%%%%%%%%%%%%%%%%%%%%%%%%%%
%%%%%%%%%%%%%%%%%%%%%%%%%%%%%%%%
\begin{frame}  \small 


\begin{table}[!htbp]
\centering
\renewcommand{\arraystretch}{1.25}
\resizebox{\textwidth}{!}{
\begin{tabular}{lccccp{6cm}}
\toprule
\textbf{Código cíclico} & \textbf{(n, k, d$_{\min}$)} & \textbf{t} & \textbf{Campo} & \textbf{Polinomio generador $g(x)$} & \textbf{Aplicaciones típicas} \\ 
\midrule
\textbf{Hamming } & (7, 4, 3) & 1 & $\mathbb{F}_2$ & $x^3 + x + 1$ & Primer código perfecto. Corrección de un solo bit en memorias ECC, microcontroladores y redes básicas. \\[4pt]\hline

\textbf{Hamming } & (15, 11, 3) & 1 & $\mathbb{F}_2$ & $x^4 + x + 1$ & Corrección de un bit. Muy usado en buses de datos y almacenamiento. \\[4pt]\hline

\textbf{BCH } & (15, 7, 5) & 2 & $\mathbb{F}_2$ & $x^8 + x^7 + x^6 + x^4 + 1$ & Corrección de dos errores. Usado en comunicaciones digitales, transmisión de datos, y en docencia como ejemplo estándar. \\[4pt]\hline

\textbf{BCH } & (31, 16, 7) & 3 & $\mathbb{F}_2$ & $x^{15} + x^{12} + x^{11} + x^{10} + x^7 + x^6 + x^4 + x^3 + 1$ & Corrección de 3 errores. Usado en radioenlaces y enlaces satelitales. \\[4pt]\hline

\textbf{Golay } & (23, 12, 7) & 3 & $\mathbb{F}_2$ & $x^{11}+x^9+x^7+x^6+x^5+x+1$ & Código perfecto binario. Usado por NASA (Voyager) y en aplicaciones de alta fiabilidad. \\[4pt]\hline

\textbf{Golay extendido } & (24, 12, 8) & 3 & $\mathbb{F}_2$ & Extensión del anterior con bit de paridad & Código perfecto extendido, usado en comunicaciones espaciales. \\[4pt]\hline


\bottomrule
\end{tabular}
}
\caption{Principales códigos cíclicos conocidos, sus parámetros y aplicaciones.}
\label{tab:cyclic_codes}
\end{table}
\end{frame} 



%%%%%%%%%%%%%%%%%%%%%%%%%%%%%%%%%%%%%%%%%%%%%%%%
%%%%%%%%%%%%%%%%%%%%%%%%%%%%%%%%%%%%%%%%%%%%%%%%
%%%%%%%%%%%%%%%%%%%%%%%%%%%%%%%%%%%%%%%%%%%%%%%%

\begin{frame}[fragile]  {generate\_test\_bmp.py}\footnotesize  

\begin{verbatim}  
usage: generate_test_bmp.py [-size]  [-out BMP_file]
\end{verbatim}
\footnotesize
\begin{verbatim}  
python3 generate_test_bmp.py --size 512 --out test512.bmp 
\end{verbatim}
\begin{center}
\includegraphics[height=0.7\textheight]{test512.pdf}
    
\end{center}

\end{frame} 


%%%%%%%%%%%%%%%%%%%%%%%%%%%%%%%%%%%%%%%%%%%%%%%%
%%%%%%%%%%%%%%%%%%%%%%%%%%%%%%%%%%%%%%%%%%%%%%%%
%%%%%%%%%%%%%%%%%%%%%%%%%%%%%%%%%%%%%%%%%%%%%%%%

\begin{frame}[fragile]  {bmp\_noisy\_channel.py}\footnotesize  

\begin{verbatim}  
usage: bmp_noisy_channel.py [-h] [--channel {bsc,ge}] [--p P] [--pg PG] 
[--pb PB] [--a A] [--b B] [--seed SEED] in_bmp out_bmp
bmp_noisy_channel.py: error: the following arguments are required: 
in_bmp, out_bmp
\end{verbatim}
\tiny
\begin{verbatim}  
python3 bmp_noisy_channel.py --channel bsc test512.bmp test512OUT_001.bmp --p 0.01
[OK] Canal: BSC(p=0.01)
     Imagen: 512x512 px, 24 bpp, stride=1536, bytes_pixeles=786432
     Bits totales: 6,291,456
     Bits alterados: 62,661  -> BER=9.959698e-03
     Bytes alterados: 60,528
[OUT] Guardado: test512OUT_001.bmp
\end{verbatim}
\begin{center}
\includegraphics[width=0.3\textwidth]{test512OUT_001.pdf}
\end{center}

\end{frame} 




%%%%%%%%%%%%%%%%%%%%%%%%%%%%%%%%%%%%%%%%%%%%%%%%

\begin{frame}[fragile]  {bmp\_noisy\_channel.py}\footnotesize  


\tiny
\begin{verbatim}  
python3 bmp_noisy_channel.py --channel bsc test512.bmp test512OUT_01.bmp --p 0.1
[OK] Canal: BSC(p=0.1)
     Imagen: 512x512 px, 24 bpp, stride=1536, bytes_pixeles=786432
     Bits totales: 6,291,456
     Bits alterados: 628,788  -> BER=9.994316e-02
     Bytes alterados: 447,473
[OUT] Guardado: test512OUT_01.bmp
\end{verbatim}
\begin{center}
\includegraphics[width=0.3\textwidth]{Tema7/test512OUT_01.pdf}
\end{center}

\end{frame} 
%%%%%%%%%%%%%%%%%%%%%%%%%%%%%%%%%%%%%%%%%%%%%%%%

\begin{frame}[fragile]  {bmp\_noisy\_channel.py}\footnotesize  


\tiny
\begin{verbatim}  
python3 bmp_noisy_channel.py --channel bsc test512.bmp test512OUT_05.bmp --p 0.5
[OK] Canal: BSC(p=0.5)
     Imagen: 512x512 px, 24 bpp, stride=1536, bytes_pixeles=786432
     Bits totales: 6,291,456
     Bits alterados: 3,145,316  -> BER=4.999345e-01
     Bytes alterados: 783,269
[OUT] Guardado: test512OUT_05.bmp
\end{verbatim}
\begin{center}

\includegraphics[width=0.3\textwidth]{Tema7/test512OUT_05.pdf}
\end{center}

\end{frame} 

%%%%%%%%%%%%%%%%%%%%%%%%%%%%%%%%%%%%%%%%%%%%%%%%

\begin{frame}[fragile]  {bmp\_noisy\_channel.py}\footnotesize  


\tiny
\begin{verbatim}  
python3 bmp_noisy_channel.py --channel bsc test512.bmp test512OUT_1.bmp --p 1  
[OK] Canal: BSC(p=1.0)
     Imagen: 512x512 px, 24 bpp, stride=1536, bytes_pixeles=786432
     Bits totales: 6,291,456
     Bits alterados: 6,291,456  -> BER=1.000000e+00
     Bytes alterados: 786,432
[OUT] Guardado: test512OUT_1.bmp
\end{verbatim}
\begin{center}
\includegraphics[width=0.3\textwidth]{Tema7/test512OUT_1.pdf}
\end{center}

\end{frame}

%%%%%%%%%%%%%%%%%%%%%%%%%%%%%%%%%%%%%%%%%%%%%%%%

\begin{frame}[fragile]  {bmp\_ecc\_channel.py (In progres...)}\tiny  
\begin{verbatim}  
usage: bmp_ecc_channel.py [-h] [--outbase OUTBASE] [--ecc {none,hamming74,bch15,golay23}] 
[--p P] [--seed SEED] [--csv CSV] in_bmp
\end{verbatim}

\begin{verbatim}  
python3 bmp_ecc_channel.py --ecc hamming74 --p 0.001  test512.bmp 
[ECC] Hamming(7,4,3)  n=7,k=4,t=1, deg(g)=4
[CANAL] BSC p=0.001  seed=123
[IMG] 512x512 px, bytes_pix=786,432, bits=6,291,456
[NO ECC] flips=6,144  BER=9.765625e-04  -> out: out_noecc.bmp
[WITH ECC] encoded_blocks=1572864, flips_on_encoded=10,961
           corrected_blocks=10,923, uncorrectable_blocks=19
[QUALITY]  exact_blocks=970,742, wrong_blocks=602,122
           post-decode bit errors vs original = 602,133  (BER=9.570646e-02)
\end{verbatim}

\begin{verbatim}  
python3 bmp_ecc_channel.py --ecc bch15 --p 0.001  test512.bmp  
[ECC] BCH(15,7,5)  n=15,k=7,t=2, deg(g)=8
[CANAL] BSC p=0.001  seed=123
[IMG] 512x512 px, bytes_pix=786,432, bits=6,291,456
[NO ECC] flips=6,144  BER=9.765625e-04  -> out: out_noecc.bmp
[WITH ECC] encoded_blocks=898780, flips_on_encoded=13,358
           corrected_blocks=13,275, uncorrectable_blocks=0
[QUALITY]  exact_blocks=898,779, wrong_blocks=1
           post-decode bit errors vs original = 1  (BER=1.589457e-07)
\end{verbatim}

\begin{verbatim}  
python3 bmp_ecc_channel.py --ecc golay23 --p 0.001  test512.bmp 
[ECC] Golay(23,12,7)  n=23,k=12,t=3, deg(g)=11
[CANAL] BSC p=0.001  seed=123
[IMG] 512x512 px, bytes_pix=786,432, bits=6,291,456
[NO ECC] flips=6,144  BER=9.765625e-04  -> out: out_noecc.bmp
[WITH ECC] encoded_blocks=524288, flips_on_encoded=12,005
           corrected_blocks=11,899, uncorrectable_blocks=0
[QUALITY]  exact_blocks=524,288, wrong_blocks=0
           post-decode bit errors vs original = 0  (BER=0.000000e+00)
\end{verbatim}

\end{frame}
%%%%%%%%%%%%%%%%%%%%%%%%%%%%%%%%
%\begin{frame}  \small  
%\begin{itemize} \item  \item \item \item 
%\end{itemize} \end{frame} 



%%%%%%%%%%%%%%%%%%%%%%%%%%%%%%%%
%%%%%%%%%%%%%%%%%%%%%%%%%%%%%%%%
%%%%%%%%%%%%%%%%%%%%%%%%%%%%%%%%
%SECCION %%%%%%%%%%%%%%%%%%%%%%%%%%%%%%%
%%%%%%%%%%%%%%%%%%%%%%%%%%%%%%%%




%%%%%%%%%%%%%%%%%%%%%%%%%%%%%%%%
%\begin{frame}  \small  
%\end{frame} 


%%%%%%%%%%%%%%%%%%%%%%%%%%%%%%%%
%\begin{frame}  \small  
%\begin{itemize} \item  \item \item \item 
%\end{itemize} \end{frame} 

%%%%%%%%%%%%%%%%%%%%%%%%%%%%%%%%
%\begin{frame}  \small \begin{ejercicio} \end{ejercicio}\end{frame} 



\begin{frame} %\frametitle{Preguntas}

%\begin{center}
%¡Gracias por su atención!\\
%\end{center}
\vspace*{2cm}

\begin{center}
\textbf{José A. Montenegro Montes} \\

\textit{Dpto. Lenguajes y Ciencias de la Computación}\\
\textit{ETSI  Informática. Universidad de Málaga}\\
\textbf{jmmontes@uma.es}\\
%\href{https://twitter.com//monteDocencia}{
%
\includegraphics[height=0.05\textheight]{twitter.jpg}}
\end{center}

\vspace*{1,5cm}


\includegraphics[height=0.15\textheight]{logoUma.jpg} \hspace*{3.6cm}

\includegraphics[height=0.15\textheight]{logoInf.jpg}   \hspace*{3.7cm}

\includegraphics[height=0.15\textheight]{logoLcc.jpg}  




\end{frame} 

\end{document} 